\cotefile{77}
\scanpage{66}
% archive sheet number: 65
\uline{Discriminant de formes quadratiques.}

Soit $I$ ens d'indices fini, $\mathscr A_I=\mathbb Z[I\times I]$ l'anneau \inserted{universelle} \deleted{\unsure{univ.}} \inserted{\unsure{le|les|des}} matrices carrées de type $I$, $\mathscr A\unsure{Q|S}_I=\mathbb Z[\operatorname{Sym}^2(I)]$ \deleted{le groupe des} \inserted{celui de la} matrices carrées symétriques universelle, \deleted{de} dont les générateurs sur $Z$ sont $\deleted{{\unsure{\mathscr A|\mathbb Z|A\mathbb Z}}^{I\times I}\longrightarrow Z}$ de la forme $\deleted{X}Y_J$, où $J$ parcourt l'ens des parties de $I$ de cardinal $1$ ou $2$. \unsure{Écriv.|Écr.|Coord.} ds $\mathscr A_I$ $\deleted{\mathscr A Q_I\longleftarrow\mathscr A Q_J}$ les $X_{ij}$\inserted{$(i,j\in I\times I)$}. \deleted{\unsure{les}} considérons
\[
\mathscr A_I\xrightarrow{\varphi}\mathscr A Q_I
\]
défini par
\[
\varphi(X_{\unsure{\{i,j|i,j|ij}})=
\begin{cases}
\deleted{X}Y_{\{i,j\}}&\text{si }i\ne j\\
2\deleted{X}Y_{\{i,j\}}=2X_{\{i\}}&\text{si }i=j.
\end{cases}
\]
qui \deleted{associ\unsure{e}} correspond à l'hom de foncteurs
\[
MS_I(k)\longrightarrow M_I(k)
\]
associant à une forme quadratique de coeff. les $a_{\{i,j\}}$, la matrice de celle-ci.

Soit $\Delta_0\in\mathscr A_I$ le déterminant de la $I$-matrice universelle, $\Delta=\varphi(\Delta_0)$ son image, qui est le discriminant de la forme quadratique universelle indexée par $I$.

On a
\[
\Delta_0=\sum_{\sigma\in\unsure{\mathfrak S|\mathscr S}_I}\varepsilon_\sigma
\overbrace{\prod_{i\in I}X_{i,\sigma i}}^{X_\sigma}
\]
Considérons la décomposition de $I$ en cycles suivant $\sigma$, le produit s'écrit -- en mettant à
\scanpage{67}
% archive sheet number: 66
part les cycles de longueur $1$ et ceux de longueur \unsure{2}
\[
\begin{aligned}
X_\sigma&=\prod_{i\in I_1(\sigma)}X_{ii}
\prod_{\deleted{\unsure{i'\in}}i\in I_2(\sigma)} X_{\{i,\sigma i\}}
\prod_{i\in I_*(\sigma)}X_{\{i,\sigma i\}}\\
\varepsilon_\sigma&=(-1)^{\unsure{\nu|v}_2+\deleted{\unsure{\nu_1}}\,\nu_*^+}
\end{aligned}
\]
où
\[
\left|\begin{aligned}
I_1(\sigma)&=I^\sigma\\
I_2(\sigma)&=\bigcup\,\text{cycles de card. }2\\
I_*(\sigma)&=\bigcup\,\text{cycles de card }\geq3\\
\nu_1&=\unsure{\text{nb}}\ \text{de cycles de card }1=\operatorname{card}\overbrace{I^\sigma}^{\inserted{I_1(\sigma)}}\deleted{\unsure{\mathrm{car}|\mathrm{card}}}\\
\nu_2&=\text{nb de cycles de card }2\\
\nu_*^+&=\text{nb de cycles de cardinal pair } >\unsure{2|4}\\
\nu_*&=\text{nb de }\rule{7em}{0.35pt}>2
\end{aligned}\right.
\]
on trouve donc
\[
\Delta=\sum_{\sigma\in\unsure{\mathfrak S|\mathscr S}_I}(-1)^{\nu_2+\nu_*^+}2^{\nu_1}
\underbrace{\prod_{i\in I_1(\sigma)}\deleted{X}Y_{\{i\}}\!
\left(\prod_{\substack{\deleted{\unsure{i\in I_2(\sigma)/\sigma}}\\H\in\dot I_2(\sigma)}}\deleted{X}Y_H\right)^{2}
\prod_{H\in\dot I_*(\sigma)}\deleted{X}Y_{H,\sigma}}_{\deleted{\unsure{X_\sigma}}\,Y_\sigma}
\]
où
\[
\deleted{X}Y_{H,\sigma}=\prod_{a\in A(H,\sigma)}X_a
\]
$A(H,\sigma)$ est l'ens des \uline{arêtes} du polygône combinatoire (orienté) dont l'ens des sommets est $H$, avec la structure polyg. orientée déduite de $\sigma|H$ \enspace $(\dot I_2(\sigma)=I_2(\sigma)/\sigma,\ \dot I_*(\sigma)=I_*(\sigma)/\sigma)$.

\deleted{Les monômes} \inserted{On peut} dire que \inserted{la donnée de} $\sigma$ \deleted{définit} \inserted{équivaut à celle d'} une \deleted{\unsure{str}} décomposition de $I$ en $I_1(\sigma)\amalg I_2(\sigma)\amalg I_3(\sigma)$ où $I_2(\sigma)$ est munie d'une involution ss pt fixe $\sigma_2$ (i.e. d'une \deleted{\unsure{de}} \unsure{rel.|rel|relation} d'équivalence à classes de cardinal $2$), \deleted{et} et $I_3(\sigma)$ d'une structure de
\scanpage{68}
% archive sheet number: 67
contour orienté, et
\[
\underset{\substack{\shortparallel\\Y_{\gamma(\sigma)}}}{Y_\sigma}
=\prod_{i\in I_1(\sigma)}Y_i
\left(\prod_{a\in I_2(\sigma)/\sigma_2}Y_a\right)^2
\prod_{a\in\underset{\substack{|\\\text{arêtes}}}{A}(I_3(\sigma))}Y_a
\]
Quant aux monômes \inserted{$Y_\sigma$}, ils correspondent $1$-$1$ aux structures de graphes \deleted{\unsure{à}} stricts \inserted{$\gamma$} sur $I$ correspondants aux $\sigma$, i.e dont les composantes connexes sont soit des sommets isolés, soit des \deleted{simplets} doublets
\figure{\begin{tikzpicture}[baseline=-.5ex]
\fill (0,0) circle (1.3pt) (.8,0) circle (1.3pt);
\draw (0,0)--(.8,0);
\end{tikzpicture}}
, soit des polygônes combinatoires d'ordre \deleted{\unsure{et|d}} $\geq3$. Pour que $Y_\sigma=Y_{\sigma'}$, il f. et s. \inserted{donc} que $\sigma$ et $\sigma'$ définissent les $=$ $I_1$ et $I_2$, et \deleted{que $I_1(\sigma)$} la $=$ structure \deleted{polyg.} de contour sur $I_3=I\setminus I_1\setminus I_2$, donc qu'ils \deleted{\unsure{le prod.|la partie}} définissent $=$ partition de $I_3$ en \deleted{parties connexes} \inserted{orbites et} que sur chaque classe $\sigma$ et $\sigma'$ sont ou bien égales, ou bien inverses. Donc chaque $Y_{\gamma\deleted{\unsure{\sigma}}}$ intervient exactement $2^{\unsure{\nu|v}_*}$ fois. Donc
\[
\begin{aligned}
\Delta={}&\sum_{\gamma\in\underset{\substack{\downarrow\\\text{ens des structures de}\\\text{graphes sur }I\text{ du type}\\\text{envisagé}}}{\Gamma_I}}
(-1)^{\nu_2(\gamma)+\nu_*^+(\gamma)}
2^{\overbrace{\nu_1(\gamma)+\nu_*(\gamma)}}\,
\underbrace{\begin{aligned}
&\prod_{i\in I_1(\unsure{\sigma|\gamma})}Y_i
\left(\prod_{a\in A(I_2(\gamma))}Y_a\right)^2\\
&\qquad\prod_{a\in A(I_3(\gamma))}Y_a
\end{aligned}}_{M_\gamma}
\end{aligned}
\]

\scanpage{69}
% archive sheet number: 68
\[
\Delta\bmod 2=\sum_{\gamma\in\Gamma_0}\deleted{\unsure{\varepsilon}}
\left(\prod_{a\in\unsure{\mathrm{Ar}(\sigma_0)|\mathrm{Ar}(\Gamma_0)|A(\gamma_0)}}Y_a\right)^2
\]
% curved pointer from the subscript of the sum
\inserted{$\hookrightarrow$ structures de graphes consistt en arêtes isolées -- i.e. \unsure{correspondt|correspondent|correspondant} aux inv. ss pts fixes de $I$ -- i.e. tels que $\unsure{\nu|v}_1(\gamma)=\nu_*(\gamma)=0$}

On retrouve que $\Delta\bmod 2=0$ si $\operatorname{card}I=n$ impair.

\deleted{Soit $B=\sum\unsure{-1}$}

Supposons pour la suite
\[
\operatorname{card}I=n=2m
\]
\deleted{Notons} et posons
\[
B=B_I=\sum_{\gamma\in\Gamma_0}\deleted{(-1)^{\unsure{m}}}
\left(\prod_{a\in\unsure{\mathrm{Ar}(I,\sigma_2)|\mathrm{Ar}(I,\gamma_2)|A(I,\gamma_2)}}Y_a\right).
\]
(\uline{NB} on aura $\deleted{n=\nu_2(\sigma)=\unsure{1}}$ $\nu_2(\gamma)=m$, $\nu_*^+(\gamma)=0$ donc $(-1)^m$ est le coeff avec lequel $\prod Y_a$ figure ds $\Delta$ \ldots) \deleted{On aura} Soit $\varepsilon=\pm1$, qu'on essaie de déterminer de façon que
\[
B^2-\varepsilon\Delta\equiv0\quad(4)
\]
i.e. tel que l'on puisse trouver $C\in\mathscr A\unsure{S|Q}_I$ avec
\[
B^2-4C=\varepsilon\Delta
\]

On \deleted{aura} \inserted{présume} que $\varepsilon=(-1)^m$.
\[
B^2=\sum_{\gamma\in\Gamma_0}\left(\prod_{a\in A(\gamma)}Y_a\right)^2
+\deleted{\unsure{2|4}}\sum_{\substack{\gamma,\gamma'\in\Gamma_0\\\gamma\ne\gamma'}}\deleted{\unsure{\varepsilon}}
\left(\prod_{a\in A(\gamma)}Y_a\right)
\left(\prod_{\unsure{a|a'}\in A(\gamma')}Y_{a'}\right)
\]

On \unsure{\deleted{f}}a
\scanpage{70}
% archive sheet number: 69
\[
\deleted{(-1)^{\unsure{n|m}}\Delta\equiv
\deleted{\sum_{\substack{\gamma\in\Gamma\\{\unsure{\nu|v}}_1(\gamma)+\nu_2}}}B+
\sum_{\gamma\in\Gamma_1}}
\]
\deleted{$\Gamma_1=$ ens des $\gamma\in\Gamma$ tels que \deleted{il y a} $\nu_1(\gamma)+\nu_*(\gamma)=1$ i.e. $\left\{\begin{aligned}
&[\nu_1(\gamma)=1,\ \nu_2(\gamma)=0\quad\text{impossible car }n\text{ pair}]\\
&\text{ou }\nu_1(\gamma)=0,\ \nu_2(\gamma)=1\quad\text{i.e pas de pt fixe\unsure{s},}\\
&\hspace{6em}\text{et seulmt une composante connexe de card }\geq3
\end{aligned}\right.$}

\deleted{Notons que $\nu_1(\gamma)+\nu_*^{\deleted{\unsure{-}}}(\gamma)\equiv m\ (2)$ puis que $m-(\nu_1(\gamma)+\nu_*(\gamma))=\nu_2(\gamma)\equiv0\ (2)$}

\deleted{Si $m$ pair, on trouve donc $\nu_1(\gamma)+\nu_*^-(\gamma)\equiv0\ (2)$ donc si $\nu_1(\gamma)+\nu_*^-(\gamma)\ne0$, on a $\unsure{\nu_1(\sigma)|\gamma_1(\sigma)|\nu_1(\gamma)}+\nu_*^-(\gamma)\geq2$ donc $2^{\nu_1(\gamma)+\nu_*(\gamma))}\equiv0\ (4)$.} \deleted{D'où} \inserted{On a} pour $m$ pair
\[
\begin{aligned}
&\deleted{\unsure{B^2}=}\\
\Delta&\equiv(-1)^m\deleted{\unsure{B^2}}\sum_{\gamma\in\Gamma_0}
\deleted{\unsure{(\varepsilon_{\sigma_2})}}\left(\prod_{a\in A(\gamma)}Y_a\right)^2\\
&\quad+2\sum_{\gamma\in\Gamma_1}(-1)^{\nu_2(\gamma)+1}
\end{aligned}
\qquad\bigg|\,\substack{\mathrm{mod}\\(4)}
\]
\margin{\fbox{\parbox{1.8cm}{$\Gamma_1$\,: graphes $\gamma\in\Gamma$ tels que $\nu_1(\gamma)=0,\ \nu_*(\gamma)=\nu_*^+(\gamma)=1$}}}
donc
\[
B^2-(-1)^m\Delta\equiv\sum_{\substack{\gamma,\gamma'\in\Gamma_0\\\gamma\ne\gamma'}}
\underbrace{\left(\prod_{a\in A(\gamma)}Y_a\right)\!\left(\prod_{a'\in A(\gamma')}Y_{a'}\right)}_{M_{\gamma,\gamma'}}
\qquad\mathrm{mod}\ 4
\]
Considérons un couple $\gamma,\gamma'$ fixé, ou ce qui revient au $=$, $2$ inv. ss pts \deleted{\unsure{fixes}} \unsure{fixes|fixées} $\sigma,\sigma'$, distinctes. \deleted{Ainsi,} Examinons la structure du monôme $M_{\gamma,\gamma'}$ : Si nous posons $I_{\sigma,\sigma'}=\{i\in I\mid\sigma i=\sigma'i\}$, on a $I=I_{\sigma,\unsure{\sigma|\sigma'}}\amalg J_{\sigma,\sigma'}$, avec $J_{\sigma,\sigma'}\ne\varnothing$. Dans
\scanpage{71}
% archive sheet number: 70
% printed form
\begin{center}
\fbox{\fbox{\shortstack{CARTES \quad DE \quad VOEUX\\POSTERS \quad -- \quad DISQUES}}}
\end{center}

\begin{center}******\end{center}

Cette année encore, l' AIDE à l'ENFANCE du VIET NAM ( dont le siège est à GIF - SUR - YVETTE -- 91190 -- 104, avenue Général Leclerc )

propose :

-- Lot de 5 cartes de vœux ( série 7 S ) \dotfill\ 10 F.

-- Lot de 7 cartes de vœux ( série IV S ) \dotfill\ 10 F.

-- Lot de 10 cartes de vœux ( série 5,6,7 ) \dotfill\ 15 F.

-- Des POSTERS au prix de \dotfill\ 20 F.

-- Des DISQUES de musique traditionnelle du Viet Nam\\
\hspace*{3em}33 tours 30 cm \dotfill\ 35 F.

-- Des rouleaux de Peinture sur Soie peints à la main . . \quad 100,00 F.

\textit{Dépositaire à l'U.S.T.L. -- Service Duplication} Monsieur HUBERSON

Poste 816

\texttt{MR \ LE PROFESSEUR GROTHENDIECK A.}\\
\texttt{MATHEMATIQUES \hspace{7em} 5}
\scanpage{72}
% archive sheet number: 71
$J_{\sigma,\sigma'}$, \deleted{\unsure{on a si}} les $a\in A(\gamma)$ et les $a'\in A(\gamma')$ sont distincts. \deleted{\unsure{Le}} le produit des $Y_a,Y_{a'}$ correspondants est aussi le monôme correspondant à une structure \deleted{polygonale} \inserted{de contour} sur $J(\sigma,\sigma')$ (dont les arêtes sont les ens de la forme $\{\unsure{i|x},\sigma i\}$ et $\{i,\sigma'i\}$), dont les composantes connexes sont d'ordre pair $\geq4$. Le couple $(\sigma,\sigma')$ est déterminé, à symétrie près \inserted{$(\sigma,\sigma')\longmapsto(\sigma',\sigma)$ sur chaque composante conn. de $J(\sigma,\sigma')$}, par la structure de graphe correspondante -- et on trouve \inserted{ainsi} toutes les structures de graphe $\gamma\in\Gamma$ pour lesquelles \deleted{il n'y a} $\unsure{\nu|v}_1(\gamma)=0$, et \deleted{les cycles} $\nu_*^+(\gamma)=\nu_*(\gamma)$. Un $=$ monôme intervient $2^{\nu_*(\gamma)}=2^{\nu_*^+(\gamma)}$ fois.

Donc on trouve
\[
\begin{aligned}
B^2-(-1)^m\Delta &\equiv
\sum_{\substack{\gamma\in\Gamma\\\boxed{\substack{\nu_1(\gamma)=0\\\nu_*(\gamma)=\nu_*^+(\gamma)>0}}}}
2^{\nu_*(\gamma)}M_\gamma\qquad\mathrm{mod}\ 4\\
&\qquad+2\sum_{\gamma\in\Gamma_1}(-1)^{n+\nu_2(\gamma)}M_\gamma\\
&\equiv2\sum_{\substack{\gamma\in\Gamma_1\\\deleted{\nu_1(\gamma)=0}\\\deleted{\nu_*(\gamma)=\nu_*^+(\gamma)=1}}}M_\gamma\qquad\not\equiv0\quad(4)!\\
&\qquad\deleted{\unsure{+}}+2\sum(-1)^{\varepsilon(\gamma)}M_\gamma\\
&=2\sum M_\gamma\bigl(\underbrace{1+(-1)^{\varepsilon(\gamma)}}_{0\ \mathrm{ou}\ 2}\bigr)\equiv0\quad(4)
\end{aligned}
\]
\scanpage{73}
% archive sheet number: 72
% printed notice
% letterhead
UNIVERSITE DES SCIENCES \hfill Montpellier, le 7 Décembre 1977\\
ET TECHNIQUES DU LANGUEDOC\\
\hspace*{2em}\rule{3.3em}{0.4pt}

Laboratoire de Chimie Physique

\begin{center}\uline{PROJECTION D'UN FILM SCIENTIFIQUE}\end{center}

Monsieur Ch. VIDAL, Professeur à l'Université de Bordeaux I présentera un film couleur de 30 minutes :

\begin{center}\textquotedbl{}LOIN DE L'EQUILIBRE\textquotedbl{}\end{center}

LE VENDREDI 16 DECEMBRE 1977 à 15 H DANS L'AMPHITHEATRE N\textdegree{} 5.

Ce film a été réalisé par le Service d'Etude de réalisation et de diffusion de documents audio-visuels (C.E.R.D.D.A.V.) du C.N.R.S. sur un thème proposé par MM. PACAULT, VIDAL, de KEPPER et Mme MERLE.

Présenté aux 2e journées internationales du film scientifique à Paris, il vient d'être projeté au C.N.R.S. le 26 Novembre au cours d'une manifestation organisée en l'honneur de I. PRIGOGINE récent prix Nobel.

\hfill J.V. ZANCHETTA.

\texttt{MR \ LE PROFESSEUR GROTHENDIECK A.}\\
\texttt{MATHEMATIQUES \hspace{7em} 5}
\scanpage{74}
% archive sheet number: 73
$A$ anneau
\[
\begin{aligned}
&\Delta\in A\\
&A_0=A/2A
\end{aligned}
\]
On veut écrire
\[
\Delta=b^2-4c\qquad b,c\in A \tag*{(1)}
\]
ce qui implique
\[
\Delta_0\in A_0^2\qquad(\Delta_0=(\Delta\bmod2))
\]
\deleted{soit}
\[
\unsure{\Delta_0|\omega_0}=b_0^2\qquad b_0\in A_0
\]
où $b_0$ est uniquement déterminé si \fbox{$A_0$ réduit}\textcircled{\scriptsize\unsure{a|*}} donc \deleted{\unsure{b}} $b$ déterminé mod $2$ \inserted{par la condition $b\equiv b_0\ (2)$}, donc si $b'$ satisfait la $=$ condition i.e $b'=b+2x$, on a
\[
{b'}^2=b^2+4(x^2+bx)\qquad\text{i.e.}\quad {b'}^{\unsure{2}}\equiv b^2\quad\text{\uuline{$(4)$}}
\]

On voit donc

\uline{Lemme} \quad $\forall\ b_0\in A_0$, il y a un unique élément \inserted{$\gamma(\deleted{\unsure{x}}b_0)$} de $A/4A$ tel que $\forall\ b\in A\deleted{\unsure{/2A}}$ avec $b\equiv b_0\ (2)$, on ait $\gamma(b_0)\equiv b^2\ (4)$. Pour $\Delta\in A$ donné, pour qu'il existe $b,c\in A$ tel qu'on ait (1), il f. et s. que

a) $\Delta_0\ (= \Delta\bmod2)$ soit le carré d'un él. $b_0$ de $A_0$ (bien déterminé car $A_0$ réduit)

b) $\gamma(b_0)=0$.

Supposons qu'on puisse trouver \deleted{un} $\unsure{v|\upsilon|u}\in A^{\deleted{\unsure{*}}}$ \deleted{\unsure{entier impair}} $\deleted{\unsure{u|v|\nu}=2\unsure{t|x|d}+1}$ \deleted{tel que les conditions précédentes sont}
\scanpage{75}
% archive sheet number: 74
% printed form
% letterhead
UNIVERSITE\\
DES SCIENCES ET TECHNIQUES \hfill Montpellier, le 1er décembre 1977\\
\uline{DU LANGUEDOC}\\
JP/MD

% slanted
\margin{\handwritten{\unsure{88|BB}}}
Diffusion : Liste A\\
n\textdegree{} 6828

AVIS DE SOUTENANCE D'UNE THESE DE DOCTORAT DE SPECIALITE

SCIENCES DE LA TERRE - Mention GEOLOGIE APPLIQUEE

\uline{Nom et prénom du candidat} : ERRE Henry

\uline{Titre de la thèse} : Contribution à l'étude de l'hydrogéologie des Corbières orientales karstiques et pseudo-karstiques, et des émergences littorales des côtes calcaires du Languedoc-Roussillon.

\uline{Date, heure et lieu de soutenance prévus} : Mardi 13 décembre 1977 à 16 h 30, salle de cours 10 (Administration).

\uline{Nom du Directeur de thèse} : Monsieur le Professeur J. AVIAS

\uline{Désignation du laboratoire où la thèse est préparée} : Centre d'Etudes et de Recherches Géologiques et Hydrologiques.

\uline{Analyse de l'ouvrage} : L'auteur étudie les terrains jura-cretacés et triasiques en tant que systèmes aquifères. Ils le conduisent à examiner plus généralement les émergences littorales des côtes calcaires et à proposer une interprétation des anomalies thermiques et chimiques observées.

\uline{Durée approximative des recherches ayant abouti à la thèse} : 2 ans 1/2

\uline{Composition du jury proposé} : M. AVIAS - Président\\
\hspace*{7em}M. BOCQUILLON\\
\hspace*{7em}M. PLEGAT\\
\hspace*{7em}M. PALOC\\
\hspace*{7em}M. SALVAYRE

Le Secrétaire Général,\\
\handwritten{\unsure{E. Siau|E. Si}}\\
\uline{E. SIAU}
\scanpage{76}
% archive sheet number: 75
\deleted{satisfaites pour} \inserted{$\unsure{v|\upsilon|u}\in A$ tel que} $\Delta v^{\deleted{\unsure{k}}2}=\deleted{\unsure{y^2-4c}}$, \deleted{\unsure{satisfasse} \illegible{2}} \deleted{(Cela signifie} \inserted{Ceci implique} que
\[
\Delta_0^{\deleted{\unsure{k}}}v_0^2=\deleted{\unsure{\Delta_0^k b_0^2}}\,{b'_0}^{2}
\]
on aura, en supposant que \textcircled{\scriptsize b} \fbox{$\Delta_0$ est él. régulier de $A_0$} que
\[
\left(\frac{b'_0}{\deleted{\unsure{\Delta_0^k}}v_0}\right)^2=\Delta_0\qquad
\text{relation dans }{A_0}_{\Delta_0}
\]
supposons que \fbox{$A_0$ soit intégralemt clos}\textcircled{\scriptsize c} alors on aura $b'_0/\deleted{\unsure{\Delta_0^k}}v_0\in A_0$, soit $b$, donc
\[
b'_0=\deleted{\unsure{\Delta^k}}b_0v_0\qquad(b_0\in A_0,\ b_0^2=\Delta_0).
\]
\deleted{(ce qui ne fait \unsure{aucun sens}}

Bien entendu, on a alors, si $b\equiv b_0\ (2)$ donc $b'=\deleted{\unsure{\Delta^k}}bv$ congru à $b'_0\ (2)$
\[
\gamma(b'_0)\equiv\underset{\substack{\shortparallel\\\deleted{\unsure{\Delta^{2k}}} b^2v^2}}{ {b'}^2}\ \deleted{\unsure{=}}(\bmod 4)
\]
et il reste à exprimer
\[
\deleted{\unsure{\Delta^{2k}}}b^2v^2\equiv\Delta^{\deleted{\unsure{2k}}+1}v^2\quad(4)
\]
on aura, si $\deleted{\unsure{\Delta v|\Delta_0}}\inserted{v\bmod4}$ est régulier dans $A/4A$. -- ce qui est le cas si \fbox{$2$ est régulier dans $A$}\textcircled{\scriptsize d} et \fbox{$v_0$ est régulier dans $A_0$} \deleted{compte tenu que $\Delta_0$ est régulier dans $A_0$}
\[
b^2\deleted{\unsure{v^2}}\equiv\Delta\quad(4)
\]
\deleted{donc}
\scanpage{77}
% printed form
% letterhead
UNIVERSITE\hfill Montpellier, le 8 Décembre 1977\\
DES SCIENCES ET TECHNIQUES\\
DU LANGUEDOC\\
------------\\
Service de la Scolarité\\
JP/GJ\\
Diffusion : Liste A\\
N°7040

\uline{AVIS DE SOUTENANCE D'UNE THESE DE DOCTORAT DE 3° CYCLE}

\uline{PHYSIQUE DU SOLIDE - APPLICATION A L' ELECTRONIQUE}

\uline{Nom et prénom du candidat} : ENSUQUE Louis

\uline{Titre de la thèse} :\textquotedbl{}Interprétation nouvelle des mécanismes de conduction dans les matériaux amorphes et désordonnés : application au B14Si\textquotedbl{}.

\uline{Date, heure et lieu de soutenance prévus} :\\
12 DECEMBRE 1977 à 10heures 30\\
Salle de réunion du C.E.E.S.

\uline{Nom du Directeur de thèse} :\\
Monsieur le Professeur PISTOULET.

\uline{Désignation du laboratoire où la thèse est préparée} :\\
Laboratoire d'Electronique C.E.E.S.

\uline{Analyse de l'ouvrage} :\\
Nous rappelons les théories actuelles concernant les corps amorphes, puis nous montrons les difficultés rencontrées pour l'interprétation de certains phénomènes physiques. Nous présentons un nouveau modèle de conduction dans les corps désordonnés et amorphes, qui tient compte de l'existence d'un désordre de composition à moyenne distance. Dans ce modèle un corps est caractérisé par la connaissance de deux quantités : énergie d'ionisation du niveau donneur ou accepteur et paramètre de désordre. Nous appliquons ce modèle au composé bore-silicium destiné à la conversion thermo électrique.

\uline{Durée approximative des recherches ayant abouti à la thèse} : 2 ans.

\uline{Composition du jury proposé} :\\
MM. : B. PISTOULET, Président\\
H. ANDRILLAT\\
J.L. ROBERT\\
J.M. DUSSEAU

Le Secrétaire Général,\\
\handwritten{\unsure{E. Siau|E. Si}}\\
\uline{E. SIAU}

% archive sheet number: 76
\scanpage{78}
Ainsi on a prouvé

\uline{Prop} Supposons $2$ régulier dans $A$, $A_0=A/2A$ réduit et \unsure{intégralement} clos, $\Delta\in A$ tel que $\Delta_0\in A_0$ soit él. régulier de $A_0$. Soit $\unsure{k=2d-|*=2d-|k=2d-1|\nu=2d+1}$, comme nous notons supra. Pour qu'il existe $b',c'\in A$, \unsure{ou on|ou un} \inserted{$\unsure{v|\upsilon|u}\in A^{\deleted{\unsure{*}}}$} a
% slanted
\margin{\uline{NB} -- \unsure{Lieu|au lieu} de $v\in A^*$, il suffit existe $v\in A$ tel que $v$ soit régulier ds $A$ et $v_0$ régulier\unsure{!?|?} ds $A_0$}
\[
\left\{\begin{aligned}
{b'}^2-4c'&=\Delta^{\deleted{\unsure{*}}}v^2\\
v_0&\text{ est $A_0$-régulier}
\end{aligned}\right.
\]
\unsure{faut s.|f. et s.} qu'il existe $b,c\in A$ tel que
\[
 b^2-4c=\Delta.
\]

\deleted{Supposons}\inserted{Considérons} maintenant \deleted{que \unsure{l'on ait|l'anneau}}\inserted{l'anneau}
\[
 A_{2\Delta}\deleted{\unsure{=A}}\inserted{=(A_2)_\Delta=(A_\Delta)_2},\ \text{localisé de $A$ pour $2\Delta$} \tag*{(\unsure{a|$\alpha$})}
\]
et l'Algèbre
\[
 B_{2\Delta}=A_{2\Delta}[T]/(T^2-\Delta)
\]
obtenue en extrayant \deleted{la} racine carrée de $\Delta$. \deleted{Suppos}\inserted{C'est une $A_{2\Delta}$-alg. étale galoisienne de grpe $\mathbb Z/2$.} \uline{Supposons} que $B_{2\Delta}$ se prolonge en une Algèbre $B_\Delta$ sur $A_\Delta$ \inserted{finie} étale \deleted{\unsure{de}} [qui sera unique : \unsure{isom|sens} unique près, car $2$ est régulier dans $A_\Delta\ldots$]. Supposons de plus \unsure{que|q.|Z}
\[
 \text{\textcircled{e}}\quad \boxed{\operatorname{Pic} A_\Delta=0}
\]
alors on sait (comme $B_\Delta$ est finie loc libre de rang $2$) que l'on peut écrire

% archive sheet number: 77
\scanpage{79}
% printed form
% letterhead
UNIVERSITE\hfill Montpellier, le 8 Décembre 1977\\
DES SCIENCES ET TECHNIQUES\\
DU LANGUEDOC\\
------------\\
Service de la Scolarité\\
JP/GJ\\
Diffusion : Liste A\\
N°7041

\uline{AVIS DE SOUTENANCE D'UNE THESE DE DOCTORAT DE}

\uline{SPECIALITE : ELECTRONIQUE DES SOLIDES}

\uline{Nom et prénom du candidat} : DAROLLES Jean-Marie

\uline{Titre de la thèse} : Etude de la conversion de l'énergie à partir de thermoéléments réfractaires Bore Silicium.

\uline{Date, heure et lieu de soutenance prévus} :\\
12 DECEMBRE 1977 à 9 Heures\\
Salle de réunion du C.E.E.S., USTL

\uline{Nom des directeurs de thèse} :\\
J.L ROBERT et J.M. DUSSEAU

\uline{Désignation du laboratoire où la thèse est préparée} :\\
Laboratoire d'Electronique C.E.E.S

\uline{Analyse de l'ouvrage} : Après avoir défini les grandeurs thermoélectriques classiques, nous montrons comment on peut optimiser le rendement d'un convertisseur. Des montages permettant d'effectuer des mesures de 77 à 1600 K sont décrits. Les résultats concernant les composés réfractaires Bore-Silicium montrent l'intérêt d'utiliser de tels corps dans la conversion thermoélectrique.

\uline{Durée approximative des recherches ayant abouti à la thèse} : 2 ans

\uline{Composition du jury proposé} :\\
MM. : B. PISTOULET, Président\\
L. LASSABATERE\\
J.-L. ROBERT\\
C. COMBESCURE\\
J.-M. DUSSEAU

Le Secrétaire Général,\\
\handwritten{\unsure{E. Siau|E. Si}}\\
\uline{E. SIAU}

% archive sheet number: 78
\scanpage{80}
\[
 B_\Delta\simeq A_\Delta[T]/(T^2+bT+c)\qquad b,c\in A_\Delta
\]
avec l'idéal discriminant de $B_\Delta$ égal :
\[
 \unsure{\mathfrak{D}|\mathscr D}(B_\Delta/A_\Delta)=(b^2-4c)A_\Delta
\]
donc, à cause de l'hypothèse $B_\Delta$ étale sur $A_\Delta$
\[
 b^2-4c\in A_\Delta^*
\]
D'ailleurs, prenant la restriction à $A_{2\Delta}$, on trouve
\[
 b^2-4c=u^2\Delta,\qquad u\in A_{2\Delta}^*=(A_\Delta)_2^*
\]

\deleted{Supposons que}
\[
 \deleted{\boxed{(A_{2\Delta})_{\unsure{2|*}}^*=\{2^\mu\Delta^\nu\unsure{v|\upsilon}\mid\mu,\nu\in\mathbb Z,\ v\in A^*)}}\ \deleted{\text{\textcircled{\unsure{e}}}}
\]
\deleted{donc on aura}
\[
 \deleted{b^2-4c=2^{2\mu}\Delta^{2\nu+1}v^2\quad\text{dans }(A_{\deleted{\unsure{2}}\Delta})_2}
\]
\deleted{avec}
\[
 \deleted{v\in A^*}
\]

Comme $b^2-4c\in A_\Delta^*$, on conclut \deleted{que}\inserted{$u^2\in A_\Delta^*$} ce qui implique
\[
 \deleted{\text{Si }\boxed{2\ \text{non inv. dans } A_\Delta}\ \text{\textcircled{f}}\ \unsure{\Longleftarrow}}
\]
$*\ u\in A_\Delta^*$ si $\boxed{A_\Delta\ \text{intégralemt clos dans } A_{2\Delta}}$ \textcircled{f}

$\deleted{\mu=0}$, donc,
\[
 b^2-4c=\Delta^{\deleted{\unsure{2\nu+1}}}u^2\qquad u\in A^* \tag*{(\unsure{*|$\alpha$})}
\]
Inversement, si on a une relation \unsure{$\alpha$|$(*)$} (en \deleted{y supposant seulement $v\in A_\Delta^*$, et oubliant les hypothèses préliminaires faites}) l'on \deleted{en aura}
\[
 \deleted{b^2-4c=\lambda^2\Delta\qquad\lambda\in A^*\quad(\lambda=(\Delta^\nu v)^{2})}
\]

% archive sheet number: 79
\scanpage{81}
% printed form
% letterhead
UNIVERSITE\\
DES SCIENCES ET TECHNIQUES\hfill Montpellier, le 1er décembre 1977\\
\uline{DU LANGUEDOC}\\
JP/MD\\
Diffusion : Liste A\\
n° 6827

AVIS DE SOUTENANCE D'UNE THESE DE DOCTORAT DE SPECIALITE

SCIENCES BIOLOGIQUES - Mention PHYSIOLOGIE ANIMALE

\uline{Nom et prénom du candidat} : FAVRE Daniel

\uline{Titre de la thèse} : Ontogenèse de l'innervation des récepteurs vestibulaires chez le chat. Etude neuro-anatomique et ultrastructurale.

\uline{Date, heure et lieu de soutenance prévus} : 10 décembre 1977 à 11 heures, salle de cours 5

\uline{Nom du Directeur de thèse} : Monsieur A. SANS

\uline{Désignation du laboratoire où la thèse est préparée} : Laboratoire de neurophysiologie sensorielle.

\uline{Analyse de l'ouvrage} : Etude de l'ontogenèse de l'innervation afférente et efférente des cellules sensorielles vestibulaires. Mise en évidence, au cours de la maturation, d'une modification structurale des contacts synaptiques entre les calices nerveux et les cellules de type I. Discussion relative à l'origine phylogénétique des cellules labyrinthiques chez les mammifères.

\uline{Durée approximative des recherches ayant abouti à la thèse} : 2 ans

\uline{Composition du jury proposé} : M. R. MARTY - Président\\
M. J. LEGRAND\\
M. R. PUJOL\\
M. A. SANS

Le Secrétaire Général,\\
\handwritten{\unsure{E. Siau|E. Si}}\\
\uline{E. SIAU}

% archive sheet number: 80
